\documentclass[tikz,border=6pt]{standalone}
\usepackage[no-math]{fontspec}
\usepackage[brazil]{babel}
\usepackage{amsmath,amssymb}
\IfFileExists{C:/Windows/Fonts/Lato_Regular.ttf}{%
  \setmainfont{Lato}[
    Path = C:/Windows/Fonts/,
    UprightFont = Lato_Regular.ttf,
    BoldFont = Lato_Bold.ttf,
    ItalicFont = Lato_Italic.ttf,
    BoldItalicFont = Lato_BoldItalic.ttf
  ]
  \setsansfont{Lato}[
    Path = C:/Windows/Fonts/,
    UprightFont = Lato_Regular.ttf,
    BoldFont = Lato_Bold.ttf,
    ItalicFont = Lato_Italic.ttf,
    BoldItalicFont = Lato_BoldItalic.ttf
  ]
}{%
  \RequirePackage[default]{lato}
}
\renewcommand{\familydefault}{\sfdefault}

\usepackage{xcolor}
\usepackage{tikz}
\usepackage{pgfplots}
\pgfplotsset{compat=1.18}
\usetikzlibrary{arrows.meta,positioning,shapes.geometric,fit,backgrounds,calc,decorations.pathreplacing,patterns}

% Paleta institucional UEFS / CHF662
\definecolor{uefsNavy}{HTML}{1B2A4A}
\definecolor{uefsTeal}{HTML}{0E7C7B}
\definecolor{uefsForest}{HTML}{2D6A4F}
\definecolor{uefsGold}{HTML}{C4983B}
\definecolor{uefsGoldDark}{HTML}{8F6B1E}
\definecolor{uefsAlert}{HTML}{B02A37}
\definecolor{uefsOrange}{HTML}{D95F02}
\definecolor{uefsInk}{HTML}{22252A}
\definecolor{uefsMuted}{HTML}{5C6470}
\definecolor{uefsWater}{HTML}{2E86AB}
\definecolor{uefsSand}{HTML}{D9B382}
\definecolor{bgNavy}{HTML}{F2F5F9}
\definecolor{bgTeal}{HTML}{E9F5F4}
\definecolor{bgForest}{HTML}{EEF5F0}
\definecolor{bgGold}{HTML}{FCF7EB}
\definecolor{bgAlert}{HTML}{FDF2F2}
\definecolor{bgGray}{HTML}{F8FAFC}
\definecolor{lineGray}{HTML}{CFD8DC}

\tikzset{
  >=Stealth,
  aaCard/.style={draw=#1, fill=white, rounded corners=4pt, line width=1pt, inner sep=6pt, align=center, font=\small},
  aaBox/.style={draw=#1, rounded corners=4pt, line width=1.1pt, inner sep=6pt, align=center, font=\footnotesize\sffamily, text=uefsInk},
  aaTitle/.style={font=\normalsize\bfseries\sffamily, text=uefsNavy, align=center},
  aaSub/.style={font=\footnotesize\sffamily, text=uefsMuted, align=center},
  aaArrow/.style={->, >=Stealth, line width=1.4pt},
  aaDArrow/.style={<->, >=Stealth, line width=1.4pt}
}


\begin{document}
\begin{tikzpicture}[
  stageBox/.style={
    draw=#1,
    fill=#1!8!white,
    rounded corners=5pt,
    line width=1.4pt,
    text width=3.15cm,
    minimum height=3.85cm,
    inner sep=7pt,
    align=center,
    font=\footnotesize\sffamily,
    text=uefsInk
  },
  stageDark/.style={
    draw=uefsAlert!55!black,
    fill=uefsAlert!75!black,
    rounded corners=5pt,
    line width=1.4pt,
    text width=3.15cm,
    minimum height=3.85cm,
    inner sep=7pt,
    align=center,
    font=\footnotesize\sffamily,
    text=white
  }
]

% Título principal (sem travessão e sem ponto de exclamação)
\node[aaTitle, font=\large\bfseries\sffamily, text=uefsNavy] (title) at (9.0, 3.05)
  {Sequência Evolutiva da Erosão Hídrica: Do Impacto da Gota à Voçoroca};

% 5 estágios com espaçamento amplo para setas nítidas
\node[stageBox=uefsTeal] (s1) at (0.0, 0.0) {%
  {\small\bfseries\color{uefsTeal}1. SPLASH}\\[1pt]
  {\footnotesize\bfseries\color{uefsTeal}(Salpicamento)}\\[6pt]
  {\color{uefsTeal!45!white}\rule{2.6cm}{0.6pt}}\\[7pt]
  Impacto da gota desagrega torrões e gera selamento%
};

\node[stageBox=uefsGoldDark] (s2) at (4.5, 0.0) {%
  {\small\bfseries\color{uefsGoldDark}2. EROSÃO}\\[1pt]
  {\footnotesize\bfseries\color{uefsGoldDark}LAMINAR}\\[6pt]
  {\color{uefsGoldDark!45!white}\rule{2.6cm}{0.6pt}}\\[7pt]
  Escoamento em lâmina remove camada fértil%
};

\node[stageBox=uefsOrange] (s3) at (9.0, 0.0) {%
  {\small\bfseries\color{uefsOrange}3. SULCOS}\\[1pt]
  {\footnotesize\bfseries\color{uefsOrange}(< 0,5 m)}\\[6pt]
  {\color{uefsOrange!45!white}\rule{2.6cm}{0.6pt}}\\[7pt]
  Fluxo concentra-se em pequenos canais (ainda reversíveis)%
};

\node[stageBox=uefsAlert] (s4) at (13.5, 0.0) {%
  {\small\bfseries\color{uefsAlert}4. RAVINAS}\\[1pt]
  {\footnotesize\bfseries\color{uefsAlert}(> 0,5 m)}\\[6pt]
  {\color{uefsAlert!45!white}\rule{2.6cm}{0.6pt}}\\[7pt]
  Incisão profunda por escoamento superficial%
};

\node[stageDark] (s5) at (18.0, 0.0) {%
  {\small\bfseries\color{white}5. VOÇOROCA}\\[1pt]
  {\footnotesize\bfseries\color{white}(Colapso)}\\[6pt]
  {\color{white!55!uefsAlert}\rule{2.6cm}{0.6pt}}\\[7pt]
  Atinge o lençol freático (piping) e solapa taludes%
};

% Setas bem visíveis entre os estágios
\draw[aaArrow, color=uefsNavy, line width=1.8pt, shorten <=3pt, shorten >=3pt] (s1.east) -- (s2.west);
\draw[aaArrow, color=uefsNavy, line width=1.8pt, shorten <=3pt, shorten >=3pt] (s2.east) -- (s3.west);
\draw[aaArrow, color=uefsNavy, line width=1.8pt, shorten <=3pt, shorten >=3pt] (s3.east) -- (s4.west);
\draw[aaArrow, color=uefsNavy, line width=1.8pt, shorten <=3pt, shorten >=3pt] (s4.east) -- (s5.west);

% Barra inferior progressiva
\draw[->, >=Stealth, line width=3.2pt, color=uefsAlert]
  ([yshift=-0.85cm]s1.south west) -- ([yshift=-0.85cm]s5.south east);

\node[font=\small\bfseries\sffamily, text=uefsAlert, align=center]
  at (9.0, -3.35)
  {Aumento exponencial da severidade, perda de solo e custo de recuperação (prevenir custa menos)};

\end{tikzpicture}
\end{document}
